\documentclass[twocolumn]{aastex631}
\usepackage{amsmath}
\usepackage{booktabs}
\shorttitle{SDSS denominator/proxy atlas}
\shortauthors{NebulaMind}
\begin{document}

\title{Supplementary SDSS Denominator and Proxy Atlas for Galaxy-Evolution Follow-up}
\author{NebulaMind Research Autopilot}
\affiliation{Public SDSS DR17 data only}

\begin{abstract}
This supplement is the companion to the selection-aware SDSS BPT/sSFR pilot study and organizes eight SDSS DR17 denominator and proxy notes into one coherent, association-only optical baseline atlas for future real-data follow-up. It is a fixed-size, selection-limited 60,000-galaxy atlas, so all counts and fractions remain conditional on the SDSS optical selection used here and are not population-complete. Because \texttt{specObjID} ordering follows SDSS targeting and plate/MJD bookkeeping, this non-random subset is not a random sky sample and introduces survey-plate and sky-coverage bias. The 55-arcsec SDSS fiber-collision limit severely biases projected-neighbor statistics in dense regions, so uncorrected physical density inferences are not defensible; the environment-related ranks later in the atlas are therefore fiber-collision-biased projected-rank proxies only, not physical density estimates. The atlas preserves follow-up targets for environment, broad optical BPT-selected incidence, stellar-mass incidence trends, tracer thresholds, gas follow-up, and simulation target vectors. Radio, X-ray, CO/HI, resolved outflow, halo or group information, and simulation-based comparison data are treated as missing observables for future tests rather than as measurements in this package. The contextual cache coverage is 24.0\% of the strict four-line S/N$\geq3$ parent, but the public parent-count cascade is not promoted to an independent retained result without query receipts. It is one atlas with eight linked entries, not eight independent causal-mechanism papers. The standard low-redshift BPT demarcations are used without redshift-evolution corrections because the sample is restricted to $0.02<z<0.12$. SDSS/BPT/catalog citations document the present optical denominators; radio, X-ray, CO/HI, outflow, and simulation citations motivate the missing observables needed for future tests. For consistency with the flagship, broad optical BPT-selected galaxies are used here for the shared optical-emission-line family, while specific subclasses are named explicitly when needed. Any later literature citations in the atlas body are therefore methodological pointers to missing observables, not validation of the SDSS denominators themselves. Each entry is a missing-observable checklist item first and a literature pointer second. \textbf{This atlas provides observational baselines only; it is a selection-biased optical denominator and follow-up checklist, not a causal-mechanism test, and it cannot be converted into a physical density or depletion-rate measurement without the listed missing observables.}
\end{abstract}

\keywords{galaxies: evolution --- surveys --- catalogs --- methods: observational --- methods: statistical}

\section{Purpose}
The main paper measures an optical association between BPT classification and catalog sSFR. These eight entries are distinct baseline-and-follow-up atlas notes: each is scientifically useful as a denominator or target-vector definition, but each lacks at least one core physical observable required by its original proposal. Although the entries span environment, maintenance heating, outflows, jet environments, mass-bin diagnostics, tracer thresholds, gas depletion, and simulation targets, they share the same optical-selection biases and missing observables. For consistency with the flagship, the atlas uses the broad optical BPT-selected family when the full optical-emission-line denominator is meant and names specific subsets only when the stricter selection matters. The BPT language and catalog-backbone language here follow the same SDSS/MPA-JHU-style value-added tables and standard demarcations as the flagship \citep{sdssdr17,brinchmann2004,york2000,baldwin1981,kewley2001,kauffmann2003bpt,kewley2006,stasinska2008,stasinska2015}. The SDSS/BPT/catalog references document the present optical denominators; the radio/X-ray/CO/HI/outflow/simulation references that appear later in the notes are role-separated as future-data motivation rather than validation of the current measurements. These are missing observables in the present catalog and are required for future mechanism tests. This is an association-only optical baseline atlas and a single internal follow-up checklist for future causal inference, not a collection of standalone mechanism papers; the limitation applies to the retained 60,000-row cache rather than to the parent SDSS DR17 release.

\section{Shared denominator limitations}
The atlas uses the same analyzed public-data backbone as the main paper: 60,000 galaxies in a fixed, selection-limited subset. The strict public four-line S/N$\geq3$ parent count of 249,917 galaxies, and the corresponding 24.0\% cache coverage, are selection-context diagnostics rather than custody-backed independent result rows; the present custody receipt inventories the 60,000-row cache and topic JSON artifacts, not the public SQL count logs behind the full parent cascade. The four-line selection is sSFR-dependent and the sample is capped and non-random, so all counts and fractions are conditional denominators rather than population-complete measurements. The galaxy-by-galaxy stellar masses and catalog sSFR values are taken from the public MPA-JHU-style \texttt{galSpecExtra} table after the same SDSS joins used in the flagship \citep{sdssdr17,brinchmann2004,york2000}. Although \texttt{PhotoObj} was joined in the catalog backbone, structural quantities such as \(R_{90}/R_{50}\), \texttt{fracDeV}, \texttt{petroR50}, and \texttt{petroR90} were not retained in the 60,000-row cache, so morphology cannot be controlled in this cycle. The SDSS/BPT/catalog references support these observed denominators; the later multiwavelength and simulation references only mark the follow-up measurements that are still missing.

The shared limitations are the same ones that appear throughout the atlas notes and should be read once here rather than repeated in each subsection. The fixed 3-arcsec SDSS fiber makes all derived optical quantities central-fiber proxies rather than global measurements, the strict four-line S/N cuts preferentially remove low-equivalent-width or passive systems, and the 55-arcsec fiber-collision limit distorts any projected-neighbor statistic in dense regions. Because the subset is sequentially selected by \texttt{specObjID}, it is also not a random sky sample and can inherit survey-plate and sky-coverage bias. These are denominator limitations, not physical results, and they apply to every subsection below. Where a subsection uses neighbor rank, the rank is an ordinal proxy inside this selection-limited sample and not a substitute for a forward-modeled physical density.

The eight subsections below are intentionally parallel baseline-plus-follow-up notes: each one states the observed optical denominator or target vector, then names the missing observables that a future multiwavelength or simulation-based test would need before any physical inference can be made. In other words, the sections are distinct follow-up domains bounded by the same optical selection bias, and their role is to organize the atlas rather than to stand as separate papers.

Public parent-count cascade values are used only as selection context in this atlas. They should be accompanied by explicit SQL/query receipts before they are promoted to independent measured results in a journal submission.

\section{Provenance map}
The candidate custody file \texttt{provenance/REAL\_DATA\_SOURCE\_CUSTODY.json} inventories the real source artifacts, byte counts, hashes, and approximate row counts without copying the source data into this candidate. Table~\ref{tab:provenance-map} maps each atlas entry and the flagship matched-control anchor to its inventoried artifact. The table is a provenance index only; it does not add measurements beyond the topic JSON files and flagship matched-pair artifact.
The custody manifest also records the row-backed source scale at a glance: the shared analysis sample is 60,000 rows, the flagship matched-pair file is 8,146 rows, and the remaining topic artifacts are summary JSONs attached to the same tracked run family.

\begin{deluxetable*}{llll}
\tabletypesize{\tiny}
\tablecaption{Custody-backed artifacts used by the atlas and flagship anchor.\label{tab:provenance-map}}
\tablehead{\colhead{Atlas entry} & \colhead{Run ID} & \colhead{Inventoried artifact} & \colhead{SHA-256}}
\startdata
Flagship matched-control anchor & \texttt{SDSS\_AGN\_SFR\_PILOT\_20260708T122000Z} & \texttt{analysis\_results.json} & \texttt{668ad7a67290600ff5028ae587d32ef239a09bd8627a480539f37e1927d659df} \\
Flagship matched-pair table & \texttt{SDSS\_AGN\_SFR\_PILOT\_20260708T122000Z} & \texttt{data/matched\_agn\_sf\_pairs.csv} & \texttt{4ea53af867cccccb2b68b81557ff84fe90ec3f13e0512ffbdc977fa7216996fd} \\
Relative neighbor-count baseline & \texttt{SDSS\_REMAINING\_TOPIC\_PILOTS\_20260708T125828Z} & \texttt{m1\_rp2\_environment\_quenching/analysis\_results.json} & \texttt{c0421620f67f3c227955affa3f4c1876cb85f8b31874d219f2cd2e35a7f9cec0} \\
Maintenance-heating denominator & \texttt{SDSS\_REMAINING\_TOPIC\_PILOTS\_20260708T125828Z} & \texttt{m1\_rp3\_maintenance\_heating/analysis\_results.json} & \texttt{06291f82c3fbe0f7fe84f7249568882ca4fa44972bcc25a55e367ef1fdcc7e6e} \\
Resolved-kinematics follow-up denominator & \texttt{SDSS\_REMAINING\_TOPIC\_PILOTS\_20260708T125828Z} & \texttt{m2\_p1\_outflow\_escape\_recycling/analysis\_results.json} & \texttt{44b2407aa691d64fd6de22eb49a8c0a185c86bb1f3b538c7bf066e904d0a3210} \\
Radio-jet environment baseline & \texttt{SDSS\_REMAINING\_TOPIC\_PILOTS\_20260708T125828Z} & \texttt{m2\_p2\_radio\_jet\_environment/analysis\_results.json} & \texttt{4e1ff701bb5b98af4945d5adad2e543e00005e1ab3907e8fae7d15e70c93e351} \\
Stellar-mass selection diagnostic & \texttt{SDSS\_REMAINING\_TOPIC\_PILOTS\_20260708T125828Z} & \texttt{m2\_p3\_feedback\_transition\_mass/analysis\_results.json} & \texttt{204ec46dc838e5e69a34b4dc2f790cb0b5e0f7fc1cb4eaa71f830779a2c92b67} \\
Tracer-threshold census & \texttt{SDSS\_REMAINING\_TOPIC\_PILOTS\_20260708T125828Z} & \texttt{m3\_p1\_multiphase\_census/analysis\_results.json} & \texttt{e711563011102657b6d5cab279c1b2ab7ed087dfc734e50932ad4edfe90d0683} \\
Low-sSFR optical denominator & \texttt{SDSS\_REMAINING\_TOPIC\_PILOTS\_20260708T125828Z} & \texttt{m3\_p2\_gas\_depletion\_efficiency/analysis\_results.json} & \texttt{42965b6f359c23b56098f5f9845561f4a4a2ba81e1e00df09fbae4acf3bcc2d9} \\
Simulation target vector & \texttt{SDSS\_REMAINING\_TOPIC\_PILOTS\_20260708T125828Z} & \texttt{m3\_p3\_simulation\_validation/analysis\_results.json} & \texttt{6f289f8c68da425eb3d8005e673bf5c5c02cf917eaa2bc6feedd053535de8f52} \\
\enddata
\tablecomments{The two flagship rows are included because the supplement repeatedly uses the main paper's 8,146-pair, -1.309 dex matched-control result as its anchor. The eight atlas-topic artifacts are under run \texttt{SDSS\_REMAINING\_TOPIC\_PILOTS\_20260708T125828Z} and use the shared 60,000-row SDSS analysis sample from \texttt{SDSS\_AGN\_SFR\_PILOT\_20260708T122000Z}. The custody manifest also records the approximate row counts for these tracked inputs, so the provenance map should be read together with the manifest when auditing scale.}
\end{deluxetable*}

\section{Atlas summary}
Table~\ref{tab:atlas-summary} condenses the follow-up menu across the eight entries. All eight entries are linked by the same limitation: they remain SDSS optical denominators or target vectors until the missing multiwavelength, morphological, or forward-model comparison data are added, so their present role is to organize follow-up rather than to establish causal physical claims.

\begin{deluxetable*}{lll}
\tabletypesize{\scriptsize}
\tablecaption{Custody-backed optical baselines in the atlas. These are selection-dependent optical baselines, not physical density, gas-mass, or feedback-efficiency measurements.\label{tab:atlas-summary}}
\tablehead{\colhead{Topic} & \colhead{Observed baseline} & \colhead{Artifact}}
\startdata
Environment & low-sSFR vs.\ 10th-neighbor rank (60,000 total; 15,000 per quartile) & \texttt{m1\_rp2} \\
Maintenance heating & broad optical BPT-selected hosts in massive low-sSFR galaxies (9,298 massive; 5,695 low-sSFR) & \texttt{m1\_rp3} \\
Outflow kinematics & high-excitation broad optical BPT-selected subset (4,440/60,000) & \texttt{m2\_p1} \\
Env.\ jets & neighbor-rank-stratified broad optical BPT-selected fraction in massive hosts & \texttt{m2\_p2} \\
Mass bin & low-sSFR and broad optical BPT-selected incidence by $M_\star$ bin (15 cells with $n\geq50$) & \texttt{m2\_p3} \\
Tracer census & tracer prevalence in 60k sample (0.136 to 0.418) & \texttt{m3\_p1} \\
Gas depletion & gas-depletion low-sSFR baseline; H$\alpha$ proxy (6,729 galaxies) & \texttt{m3\_p2} \\
Simulation vector & mass-redshift target vector (15 cells with $n\geq50$) & \texttt{m3\_p3} \\
\enddata
\tablecomments{The table is a compact index of the subsection-level missing-observables lists; it does not add new measurements or change any counts. The sharp retention drop at higher S/N mainly reflects the optical emission-line selection function, which preferentially removes low-equivalent-width or passive systems from the denominator; the surviving sample therefore becomes less representative of quiescent hosts as the cut tightens.}
\end{deluxetable*}

\begin{deluxetable*}{lll}
\tabletypesize{\scriptsize}
\tablecaption{Required follow-up data for interpreting the optical baselines in Table~\ref{tab:atlas-summary}.\label{tab:atlas-followup}}
\tablehead{\colhead{Topic} & \colhead{Missing observables} & \colhead{Future follow-up domain}}
\startdata
Environment & group catalogs; central/satellite labels; halo mass; fiber-collision correction & environment test \\
Maintenance heating & X-ray cavities; cooling luminosity; radio jet powers; halo-selected parents & radio/X-ray follow-up \\
Outflow kinematics & resolved velocities; halo potentials; multiphase gas; CGM tracers & kinematic follow-up \\
Env.\ jets & radio morphology/age; cavity energetics; hot-gas density & radio-jet follow-up \\
Mass bin & gas fractions; baryon deficits; halo masses; physical-observable constraints & selection diagnostic \\
Tracer census & multiphase tracers; shared denominator; aperture model & multiphase follow-up \\
Gas depletion & CO/dust gas masses; aperture-matched SFRs; morphology; environment & CO/HI follow-up \\
Simulation vector & simulations through SDSS/MaNGA/ALMA/X-ray/radio selection; aperture/noise models & forward model \\
\enddata
\tablecomments{These rows are missing-observable requirements, not measurements in the current SDSS-only atlas.}
\end{deluxetable*}

\section{Atlas notes}
As a reminder, each atlas entry is a baseline-plus-follow-up checklist, not a standalone physical-mechanism result.

\subsection{Relative neighbor-count baseline: SDSS 10th-neighbor index for low-sSFR incidence}
Artifact \texttt{m1\_rp2\_environment\_quenching/analysis\_results.json} establishes a relative neighbor-count baseline within the emission-line denominator that can later be joined to group catalogs and halo masses. The SDSS emission-line denominator contains 60,000 galaxies with an internally computed 10th-neighbor index. The high-index quartile has a low-sSFR emission-line fraction of 0.230 (3,456/15,000), while the low-index quartile has 0.181 (2,710/15,000). The bootstrap high-minus-low interval is [0.041, 0.059], and a linear probability model adjusted for log stellar mass and redshift gives a high-index coefficient of 0.032 +/- 0.004, corresponding to an approximate 3.2 percentage-point increase in low-sSFR incidence at fixed mass and redshift.
\par\noindent\textbf{Fiber-collision warning.} The 10th-neighbor index is the rank of the 10th nearest companion in projected sky separation within the full $0.02<z<0.12$ slice, with no additional line-of-sight velocity window. It is an internal ordinal rank within this selection-biased sample and should not be interpreted as a physical environmental volume density or halo density. The SDSS 55-arcsec fiber-collision limit systematically removes close spectroscopic neighbors in dense regions before any physical interpretation is attempted; this is the angular collision scale defined in the SDSS targeting paper \citep{strauss2002}, and at the sample redshifts it corresponds to a meaningful projected pair separation. A single transverse distance at the sample median redshift is not quoted because the present custody inventory does not include a median-redshift result receipt for this atlas note; the safe statement is that the collision scale is astrophysically relevant across the stated redshift interval. Standard mitigation would require explicit fiber-collision corrections, nearest-neighbor redshift assignment or angular-correlation corrections, overlapping-plate recovery where available, and preferably a group/halo catalog with central/satellite labels; SDSS correlation-function and halo-catalog work demonstrates why those corrections are separate data products rather than optional prose caveats \citep{zehavi2002,zehavi2011,blanton2003,guo2012,yang2007,tinker2021}. Compared with SDSS group catalogs such as the Yang et al.\ halo-based catalog and later probabilistic group-catalog methods, this 10th-neighbor index lacks halo mass, central/satellite classification, and membership probabilities, so it should be read only as a fiber-collision-biased projected-neighbor rank. The follow-up ingredients are group catalogues, robust central/satellite labels, halo masses, a spectroscopic fiber-collision correction at the 55-arcsec scale, morphology, and multi-redshift selection functions. This entry remains an optical baseline only; the missing observables listed in Table~\ref{tab:atlas-followup} are required before any physical inference.
\begin{figure}
\centering
\includegraphics[width=\columnwidth]{../figures/topic-01.pdf}
\caption{SDSS optical emission-line denominator: the low-sSFR emission-line fraction as a function of the 10th-neighbor index in the SDSS emission-line sample. This is a selection-dependent baseline for future group- and halo-matched follow-up.}
\label{fig:m1-rp2-neighbor-count-baseline}
\end{figure}


\subsection{Maintenance-heating denominator: broad optical BPT-selected hosts in massive SDSS galaxies}
Artifact \texttt{m1\_rp3\_maintenance\_heating/analysis\_results.json} isolates the broad optical BPT-selected duty-cycle denominator that radio and X-ray data would need to test maintenance heating. Among massive, low-sSFR SDSS emission-line galaxies, the broad optical BPT-selected fraction can serve as a denominator for X-ray and radio maintenance-heating follow-up. The massive subset (\(\log M_\star \geq 10.8\)) contains 9,298 emission-line galaxies, of which 5,695 are low-sSFR by the pilot threshold. The broad optical BPT-selected fraction is 0.430 in the massive subset and 0.607 among massive low-sSFR objects. This provides an optical duty-cycle denominator for X-ray and radio follow-up, not a heating-to-cooling measurement. Future physical validation requires X-ray cavity or cooling-luminosity measurements \citep{fabian2012}, calibrated radio jet mechanical powers \citep[e.g.,][]{best2005,hardcastle2020}, halo-selected parent catalogues, and nondetection modelling. Optical broad BPT selection primarily traces the radiative-mode denominator, so it cannot isolate the mechanically dominated jet-mode population without contemporaneous X-ray and radio measurements \citep{heckmanbest2014}. This entry remains an optical baseline only; the missing observables listed in Table~\ref{tab:atlas-followup} are required before any physical inference.

\begin{figure}
\centering
\includegraphics[width=\columnwidth]{../figures/topic-02.pdf}
\caption{SDSS optical emission-line denominator: the massive and low-sSFR SDSS emission-line subsets used as a baseline for future X-ray and radio measurements.}
\label{fig:m1-rp3-maintenance-heating}
\end{figure}


\subsection{High-excitation broad optical BPT-selected baseline: resolved kinematics follow-up}
Artifact \texttt{m2\_p1\_outflow\_escape\_recycling/analysis\_results.json} isolates the high-excitation broad optical BPT-selected denominator that resolved kinematics would need to test escape versus recycling. High-excitation broad optical BPT-selected candidates number 4,440 of 60,000 emission-line galaxies (0.074). Their median \(\log {\rm sSFR}\) is -11.53, compared with -10.14 for the full denominator. SDSS does not measure escape velocity or multiphase outflow velocities here; the note supplies a denominator for resolved follow-up rather than an escape or recycling result. The follow-up ingredients are resolved outflow velocities, halo potentials, molecular, ionized, and neutral gas phases, and CGM recycling tracers. Without IFU kinematics to decouple non-circular outflow components from host rotation, optical excitation alone cannot determine whether the gas exceeds the halo escape speed \citep{harrison2018}. This entry remains an optical baseline only; the missing observables listed in Table~\ref{tab:atlas-followup} are required before any physical inference.

\begin{figure}
\centering
\includegraphics[width=\columnwidth]{../figures/topic-03.pdf}
\caption{SDSS optical emission-line denominator: the high-excitation BPT-selected subset used to define an observational baseline for future resolved-kinematic measurements.}
\label{fig:m2-p1-outflow-escape-recycling}
\end{figure}


\subsection{Radio-jet environment baseline: broad optical BPT-selected fraction vs. 10th-neighbor index in massive hosts}
Artifact \texttt{m2\_p2\_radio\_jet\_environment/analysis\_results.json} defines the environment-stratified optical denominator that future radio and X-ray work could test. This subsection reuses the same projected-neighbor ranking described in the relative neighbor-count baseline above and motivates environment-stratified radio and X-ray follow-up. Among massive hosts, the high-index quartile has a broad optical BPT-selected fraction of 0.509, while the low-index quartile has 0.367. The bootstrap high-minus-low interval is [0.112, 0.170]. This is an optical/environment denominator for future radio-jet follow-up; it does not measure radio jet power or coupling efficiency. The follow-up ingredients are radio jet morphology and age, cavity or shock energetics, hot-gas density, and calibrated jet-power estimates. This entry remains an optical baseline only; the missing observables listed in Table~\ref{tab:atlas-followup} are required before any physical inference.

\begin{figure}
\centering
\includegraphics[width=\columnwidth]{../figures/topic-04.pdf}
\caption{SDSS optical emission-line denominator: the high- and low-index quartile comparison among massive SDSS hosts, used as a baseline for future radio-jet and X-ray work.}
\label{fig:m2-p2-radio-jet-environment}
\end{figure}


\subsection{Stellar-mass selection diagnostic: low-sSFR and broad optical BPT-selected incidence}
Artifact \texttt{m2\_p3\_feedback\_transition\_mass/analysis\_results.json} measures the incidence of low catalog-sSFR and broad optical BPT-selected classification across stellar-mass bins in this optical-emission-line denominator. The first stellar-mass bin with low-sSFR fraction above 0.5 is \(\log(M_\star/M_\odot) \in [11.0,12.5]\), and the broad optical BPT-selected incidence peaks in the 11.0--12.5 bin at 0.520 within this selection-limited, SpecObjID-capped pilot sample. In this optical-emission-line denominator, that peak is consistent with a selection-function bias: the S/N$\geq$3 cut preferentially removes truly passive, massive galaxies, and at the highest redshifts the same flux-limited cut is most severe for lower-mass galaxies, so the surviving mass distribution is not completeness-corrected across the full bin range. It must not be interpreted as a universal physical threshold. This is an optical distribution diagnostic; gas fractions and baryon deficits are needed before assigning any physical meaning to the apparent incidence change. The follow-up ingredients are gas fractions, baryon deficits, halo masses, central velocity dispersion proxies \citep{piotrowska2022}, stellar-regulation observables, and high-redshift extensions. The same binning is therefore best treated as a population-distribution diagnostic, not a physical transition mass for individual galaxies \citep{peng2010,wetzel2013,dekel2006}. This entry remains an optical baseline only; the missing observables listed in Table~\ref{tab:atlas-followup} are required before any physical inference.

\begin{figure}
\centering
\includegraphics[width=\columnwidth]{../figures/topic-05.pdf}
\caption{SDSS optical emission-line denominator: mass-bin diagnostic for low-sSFR and broad optical BPT-selected incidence in the SDSS emission-line denominator. This is a population baseline for future gas-inclusive follow-up.}
\label{fig:m2-p3-transition-mass}
\end{figure}


\subsection{Tracer-threshold census for multiphase follow-up}
Artifact \texttt{m3\_p1\_multiphase\_census/analysis\_results.json} compares optical tracer choices against one shared denominator before any multiphase census is attempted. Simple optical tracer definitions change the inferred broad optical BPT-selected prevalence within one common SDSS denominator. Within the same 60,000-galaxy denominator, simple optical tracer definitions produce prevalence from 0.136 to 0.418. The widest-to-narrowest prevalence ratio is 3.1 before adding molecular, neutral, X-ray, or radio phases. This demonstrates why a common-denominator multiphase census is required; it does not measure molecular or neutral outflow rates. The follow-up ingredients are ionized, molecular, and neutral tracers, X-ray or radio tracers, a shared parent denominator, and a consistent aperture model. This entry remains an optical baseline only; the missing observables listed in Table~\ref{tab:atlas-followup} are required before any physical inference.

\begin{figure}
\centering
\includegraphics[width=\columnwidth]{../figures/topic-06.pdf}
\caption{SDSS optical emission-line denominator: prevalence of alternative tracer definitions within the 60,000-galaxy sample. This is a baseline for future multiphase work.}
\label{fig:m3-p1-multiphase-census}
\end{figure}


\subsection{Low-sSFR optical denominator: baseline for future CO/HI gas measurements}
Artifact \texttt{m3\_p2\_gas\_depletion\_efficiency/analysis\_results.json} defines the denominator for CO/HI gas-fraction and depletion-time follow-up. Using the gas-depletion note's low-sSFR baseline, the massive low-sSFR denominator contains 6,729 galaxies in the SDSS emission-line sample. This denominator is note-specific and should not be conflated with the \(\log M_\star \geq 10.8\) maintenance-heating subset summarized above. Its broad optical BPT-selected fraction is 0.549, and the median H-alpha luminosity proxy is \(\log (L_{\mathrm{H}\alpha}/\mathrm{erg\,s^{-1}}) = 40.061\). Here the H-alpha luminosity proxy is the aperture-corrected \texttt{galSpecExtra} catalog value from the retained result artifact rather than raw fiber flux or a directly observed total gas quantity; that catalog-level correction extrapolates the fiber measurement beyond the aperture in a model-dependent way and assumes line emission broadly tracks the broadband light profile. The median H-alpha luminosity proxy is 0.66 dex lower than in massive star-forming emission-line galaxies. The stellar masses and sSFR values remain on the MPA-JHU/Brinchmann et al.\ Kroupa-IMF catalog scale used by \texttt{galSpecExtra}; no IMF conversion is applied here, so future \(M_{\rm gas}/M_\star\) comparisons should either keep that scale or state an explicit conversion. SDSS optical data alone cannot distinguish bulk molecular-gas depletion from localized reductions in star-formation efficiency or measure total cold-gas mass \citep{tacconi2018}; this note identifies the CO/HI follow-up denominator and optical baseline required for spatially resolved gas tests. Future physical validation requires global CO or dust-based molecular gas masses \citep[e.g., xCOLD GASS;][]{xcoldgass2017}, atomic hydrogen masses \citep[e.g., xGASS;][]{xgass2018}, aperture-matched SFRs, morphology, and environment labels, together with radio, X-ray, and IFU observables that are still missing here. Existing MaNGA+ALMA survey designs provide a practical blueprint for the needed aperture-matched stellar-population, ionized-gas, and CO measurements, but no such cross-match is performed in this atlas \citep{lin2020}. Any conversion from CO luminosity to molecular-gas mass would also require an explicit CO-to-H$_2$ conversion factor and metallicity/excitation assumptions \citep{bolatto2013,sandstrom2013}; no such conversion is performed in this optical atlas. The MPA-JHU H$\alpha$-based proxy is dust-corrected at the catalog level with a Balmer-decrement attenuation prescription consistent with Charlot \& Fall (2000), so the quantity here is model-dependent even before any later gas interpretation \citep{charlotfall2000}. In BPT-selected hosts, non-stellar ionization, Balmer absorption, and line-ratio-dependent attenuation can make standard H$\alpha$/H$\beta$ dust conversions less stable than in purely star-forming regions, so the catalog proxy should not be recast as an aperture-matched SFR without an explicit AGN-host attenuation model. As with any H$\alpha$-based proxy, residual dust attenuation and stellar-absorption systematics can still affect the optical denominator, so this value should be read as a line-luminosity proxy rather than a direct total cold-gas-mass measurement. This entry remains an optical baseline only; the missing observables listed in Table~\ref{tab:atlas-followup} are required before any physical inference. No mock, synthetic, fake, placeholder, or toy data were used in compiling this atlas.

\begin{figure}
\centering
\includegraphics[width=\columnwidth]{../figures/topic-07.pdf}
\caption{SDSS optical emission-line denominator: the massive low-sSFR SDSS galaxies available for CO/HI depletion-time follow-up.}
\label{fig:m3-p2-gas-depletion-efficiency}
\end{figure}


\subsection{Simulation target vector for forward-model comparison}
Artifact \texttt{m3\_p3\_simulation\_validation/analysis\_results.json} provides a compact observed target vector for forward modelling. The pilot writes 15 mass-redshift cells with \(n \geq 50\) as a compact comparison vector for low-sSFR fraction, broad optical BPT-selected incidence, and median \(u-r\) colour versus mass and redshift. Across mass bins, low-sSFR fractions span 0.005-0.729, and broad optical BPT-selected fractions span 0.003-0.520.

\begin{deluxetable*}{lrrrrr}
\tabletypesize{\tiny}
\tablecaption{Mass-redshift target vector retained in \texttt{m3\_p3\_simulation\_validation/analysis\_results.json}.\label{tab:simulation-vector}}
\tablehead{\colhead{$\log(M_\star/M_\odot)$ bin} & \colhead{$z$ bin} & \colhead{$N$} & \colhead{Low-sSFR fraction} & \colhead{Broad optical BPT fraction} & \colhead{Median \(u-r\)}}
\startdata
8.0--9.5 & 0.02--0.05 & 6,201 & 0.006 & 0.003 & 1.532 \\
8.0--9.5 & 0.05--0.08 & 1,638 & 0.001 & 0.001 & 1.379 \\
8.0--9.5 & 0.08--0.12 & 300 & 0.007 & 0.010 & 1.045 \\
9.5--10.0 & 0.02--0.05 & 3,607 & 0.061 & 0.030 & 1.854 \\
9.5--10.0 & 0.05--0.08 & 6,059 & 0.013 & 0.008 & 1.696 \\
9.5--10.0 & 0.08--0.12 & 2,187 & 0.003 & 0.001 & 1.516 \\
10.0--10.5 & 0.02--0.05 & 2,962 & 0.256 & 0.154 & 2.264 \\
10.0--10.5 & 0.05--0.08 & 7,581 & 0.161 & 0.090 & 2.119 \\
10.0--10.5 & 0.08--0.12 & 8,593 & 0.062 & 0.040 & 1.920 \\
10.5--11.0 & 0.02--0.05 & 1,895 & 0.581 & 0.430 & 2.623 \\
10.5--11.0 & 0.05--0.08 & 5,083 & 0.451 & 0.297 & 2.580 \\
10.5--11.0 & 0.08--0.12 & 9,861 & 0.326 & 0.209 & 2.455 \\
11.0--12.5 & 0.02--0.05 & 390 & 0.856 & 0.610 & 2.830 \\
11.0--12.5 & 0.05--0.08 & 1,199 & 0.805 & 0.563 & 2.851 \\
11.0--12.5 & 0.08--0.12 & 2,444 & 0.672 & 0.485 & 2.838 \\
\enddata
\tablecomments{Fractions and median \(u-r\) values are rounded to three decimals for compact display; the retained JSON stores the full precision. Each low-sSFR and broad optical BPT fraction is computed directly from the row-specific \(N\) count, i.e. a bin count divided by the bin denominator. This table is an observed optical target vector, not a simulation comparison.}
\end{deluxetable*}

The follow-up ingredients are simulations \citep{eagle2015} passed through the exact optical S/N and fiber-aperture selection function used here, including the sequential 60,000-row \texttt{specObjID} cache cap, then through the SDSS, MaNGA, ALMA, X-ray, and radio selection functions, together with aperture models and noise models. Forward modelling should also synthesize optical emission-line classifications, including BPT selection, before comparison to the atlas vector \citep{hirschmann2017}. Without those matched selection steps, any simulation comparison is not a valid test. This entry remains an optical baseline only; the missing observables listed in Table~\ref{tab:atlas-followup} are required before any physical inference.

\begin{figure}
\centering
\includegraphics[width=\columnwidth]{../figures/topic-08.pdf}
\caption{SDSS optical emission-line denominator: low-sSFR fraction, broad optical BPT-selected incidence, and colour versus mass and redshift in the SDSS emission-line sample. This is an observed target vector for forward modelling.}
\label{fig:m3-p3-simulation-validation}
\end{figure}

\section{Package decision}
These eight entries should remain supplementary until the missing observables are added. They are best treated as a single internal follow-up checklist: suitable as follow-up target definitions, denominator baselines, or appendix material under the main result, but not as independent causal-mechanism papers in their current SDSS-only form. The observational role of each subsection is therefore to define what real data are still needed, not to expand the current SDSS-only measurement set. In particular, this catalog-proxy atlas does not convert optical denominators into physical densities, depletion rates, or causal feedback claims.

\section*{Data Availability}
This atlas uses public SDSS DR17 spectroscopy, photometry, emission-line measurements, and MPA-JHU-style value-added catalog tables only. No proprietary data were used. The custody receipt is \texttt{provenance/REAL\_DATA\_SOURCE\_CUSTODY.json}. It inventories the shared \texttt{SDSS\_AGN\_SFR\_PILOT\_20260708T122000Z/data/analysis\_sample\_bpt.csv} source cache (60,000 rows; SHA-256 \texttt{6f982fa5778c3900239149b28729f701390fe393a164b95236229adc1e422883}), the flagship matched-control anchor \texttt{analysis\_results.json} (SHA-256 \texttt{668ad7a67290600ff5028ae587d32ef239a09bd8627a480539f37e1927d659df}) and \texttt{data/matched\_agn\_sf\_pairs.csv} (8,146 rows; SHA-256 \texttt{4ea53af867cccccb2b68b81557ff84fe90ec3f13e0512ffbdc977fa7216996fd}), plus the eight atlas JSON artifacts listed in Table~\ref{tab:provenance-map}. All eight notes remain conditional on the optical-selection denominators summarized in this atlas. This candidate package does not contain the executable analysis scripts or frozen command recipes that generated the downstream JSON files, so full regeneration from the candidate alone is not claimed. No mock, synthetic, fake, placeholder, or toy data were used.

\facilities{SDSS}

\begin{thebibliography}{}
\bibitem[Abdurro'uf et al.(2022)]{sdssdr17} Abdurro'uf, Accetta, K., Aerts, C., et al. 2022, ApJS, 259, 35; source identifier unverified / do not integrate
\bibitem[Baldwin et al.(1981)]{baldwin1981} Baldwin, J.~A., Phillips, M.~M., \& Terlevich, R. 1981, PASP, 93, 5; source identifier unverified / do not integrate
\bibitem[Best et al.(2005)]{best2005} Best, P.~N., Kauffmann, G., Heckman, T.~M., et al. 2005, MNRAS, 362, 25; ADS bibcode: 2005MNRAS.362...25B; DOI: 10.1111/j.1365-2966.2005.09192.x
\bibitem[Blanton et al.(2003)]{blanton2003} Blanton, M.~R., Lin, H., Lupton, R.~H., et al. 2003, ApJ, 592, 819; source identifier unverified / do not integrate
\bibitem[Bolatto et al.(2013)]{bolatto2013} Bolatto, A.~D., Wolfire, M., \& Leroy, A.~K. 2013, ARA\&A, 51, 207; DOI: 10.1146/annurev-astro-082812-140944
\bibitem[Brinchmann et al.(2004)]{brinchmann2004} Brinchmann, J., Charlot, S., White, S.~D.~M., et al. 2004, MNRAS, 351, 1151; source identifier unverified / do not integrate
\bibitem[Cid Fernandes et al.(2011)]{cidfernandes2011} Cid Fernandes, R., Stasi{\'n}ska, G., Schlickmann, S., et al. 2011, MNRAS, 413, 1687; source identifier unverified / do not integrate
\bibitem[Hardcastle \& Croston(2020)]{hardcastle2020} Hardcastle, M.~J., \& Croston, J.~H. 2020, New Astronomy Reviews, 88, 101539; arXiv: 2003.06137
\bibitem[Harrison et al.(2018)]{harrison2018} Harrison, C.~M., Costa, T., Tadhunter, C.~N., et al. 2018, Nature Astronomy, 2, 198; ADS bibcode: 2018NatAs...2..198H
\bibitem[Saintonge et al.(2017)]{xcoldgass2017} Saintonge, A., Catinella, B., Tacconi, L.~J., et al. 2017, ApJS, 233, 22; ADS bibcode: 2017ApJS..233...22S
\bibitem[Catinella et al.(2018)]{xgass2018} Catinella, B., Saintonge, A., Janowiecki, S., et al. 2018, MNRAS, 476, 875; ADS bibcode: 2018MNRAS.476..875C
\bibitem[Stasi{\'n}ska et al.(2008)]{stasinska2008} Stasi{\'n}ska, G., Asari, N.~V., Cid Fernandes, R., et al. 2008, MNRAS, 391, L29; source identifier unverified / do not integrate
\bibitem[Stasi{\'n}ska et al.(2015)]{stasinska2015} Stasi{\'n}ska, G., Costa Duarte, M.~V., Vale Asari, N., Cid Fernandes, R., \& Sodr{\'e}, L. 2015, MNRAS, 449, 559; source identifier unverified / do not integrate
\bibitem[Strauss et al.(2002)]{strauss2002} Strauss, M.~A., Weinberg, D.~H., Lupton, R.~H., et al. 2002, AJ, 124, 1810; ADS bibcode: 2002AJ....124.1810S; DOI: 10.1086/342343
\bibitem[Dekel \& Birnboim(2006)]{dekel2006} Dekel, A., \& Birnboim, Y. 2006, MNRAS, 368, 2; source identifier unverified / do not integrate
\bibitem[Fabian(2012)]{fabian2012} Fabian, A.~C. 2012, ARA\&A, 50, 455; source identifier unverified / do not integrate
\bibitem[Heckman \& Best(2014)]{heckmanbest2014} Heckman, T.~M., \& Best, P.~N. 2014, ARA\&A, 52, 589; source identifier unverified / do not integrate
\bibitem[Guo et al.(2012)]{guo2012} Guo, H., Zehavi, I., Zheng, Z., et al. 2012, MNRAS, 427, 428; source identifier unverified / do not integrate
\bibitem[Hirschmann et al.(2017)]{hirschmann2017} Hirschmann, M., Charlot, S., Feltre, A., et al. 2017, MNRAS, 472, 115; DOI: 10.1093/mnras/stx1907
\bibitem[Kauffmann et al.(2003a)]{kauffmann2003bpt} Kauffmann, G., Heckman, T.~M., Tremonti, C., et al. 2003a, MNRAS, 346, 1055; ADS bibcode: 2003MNRAS.346.1055K
\bibitem[Kewley et al.(2001)]{kewley2001} Kewley, L.~J., Dopita, M.~A., Sutherland, R.~S., Heisler, C.~A., \& Trevena, J. 2001, ApJ, 556, 121; source identifier unverified / do not integrate
\bibitem[Kewley et al.(2006)]{kewley2006} Kewley, L.~J., Groves, B., Kauffmann, G., \& Heckman, T. 2006, MNRAS, 372, 961; source identifier unverified / do not integrate
\bibitem[McNamara \& Nulsen(2007)]{mcnamara2007} McNamara, B.~R., \& Nulsen, P.~E.~J. 2007, ARA\&A, 45, 117; source identifier unverified / do not integrate
\bibitem[Lin et al.(2020)]{lin2020} Lin, L., Ellison, S.~L., Pan, H.-A., et al. 2020, ApJ, 903, 150; DOI: 10.3847/1538-4357/abba69
\bibitem[Peng et al.(2010)]{peng2010} Peng, Y.-j., Lilly, S.~J., Kovac, K., et al. 2010, ApJ, 721, 193; source identifier unverified / do not integrate
\bibitem[Piotrowska et al.(2022)]{piotrowska2022} Piotrowska, J.~M., Bluck, A.~F.~L., Maiolino, R., \& Peng, Y.-j. 2022, MNRAS, 512, 1052; ADS bibcode: 2022MNRAS.512.1052P; DOI: 10.1093/mnras/stac532
\bibitem[Sandstrom et al.(2013)]{sandstrom2013} Sandstrom, K.~M., Leroy, A.~K., Walter, F., et al. 2013, ApJ, 777, 5; ADS bibcode: 2013ApJ...777....5S
\bibitem[Charlot \& Fall(2000)]{charlotfall2000} Charlot, S., \& Fall, S.~M. 2000, ApJ, 539, 718; ADS bibcode: 2000ApJ...539..718C
\bibitem[Schaye et al.(2015)]{eagle2015} Schaye, J., Crain, R.~A., Bower, R.~G., et al. 2015, MNRAS, 446, 521; source identifier unverified / do not integrate
\bibitem[Tacconi et al.(2018)]{tacconi2018} Tacconi, L.~J., Genzel, R., Saintonge, A., et al. 2018, ApJ, 853, 179; ADS bibcode: 2018ApJ...853..179T
\bibitem[Tinker(2021)]{tinker2021} Tinker, J.~L. 2021, ApJ, 911, 52; DOI: 10.3847/1538-4357/abe044
\bibitem[Wetzel et al.(2013)]{wetzel2013} Wetzel, A.~R., Tinker, J.~L., Conroy, C., \& van den Bosch, F.~C. 2013, MNRAS, 432, 336; source identifier unverified / do not integrate
\bibitem[Yang et al.(2007)]{yang2007} Yang, X., Mo, H.~J., van den Bosch, F.~C., et al. 2007, ApJ, 671, 153; ADS bibcode: 2007ApJ...671..153Y
\bibitem[York et al.(2000)]{york2000} York, D.~G., Adelman, J., Anderson, J.~E., Jr., et al. 2000, AJ, 120, 1579; source identifier unverified / do not integrate
\bibitem[Zehavi et al.(2002)]{zehavi2002} Zehavi, I., Blanton, M.~R., Frieman, J.~A., et al. 2002, ApJ, 571, 172; ADS bibcode: 2002ApJ...571..172Z
\bibitem[Zehavi et al.(2011)]{zehavi2011} Zehavi, I., Zheng, Z., Weinberg, D.~H., et al. 2011, ApJ, 736, 59; DOI: 10.1088/0004-637X/736/1/59
\end{thebibliography}

\end{document}
